\documentclass[11pt,a4paper]{article}
\usepackage[top=2.3cm,bottom=2.3cm,left=2.2cm,right=2.2cm]{geometry}

\usepackage{amsmath,amssymb}
\usepackage{fontspec}
\setmainfont{Liberation Serif}
\setsansfont{Liberation Sans}
\setmonofont{DejaVu Sans Mono}[Scale=0.84]

\usepackage{xcolor}
\definecolor{brandblue}{RGB}{20, 55, 105}
\definecolor{accentblue}{RGB}{35, 95, 165}
\definecolor{codebg}{RGB}{248, 249, 251}
\definecolor{codeframe}{RGB}{215, 222, 230}
\definecolor{javakeyword}{RGB}{0, 45, 170}
\definecolor{javastring}{RGB}{5, 120, 20}
\definecolor{javacomment}{RGB}{120, 120, 120}
\definecolor{javanumber}{RGB}{20, 75, 220}

\definecolor{noteborder}{RGB}{30, 110, 65}
\definecolor{notebg}{RGB}{242, 249, 245}
\definecolor{warnborder}{RGB}{185, 75, 10}
\definecolor{warnbg}{RGB}{254, 248, 240}
\definecolor{tipborder}{RGB}{30, 95, 160}
\definecolor{tipbg}{RGB}{242, 247, 254}

\usepackage{fancyhdr}
\usepackage{lastpage}
\pagestyle{fancy}
\fancyhf{}
\lhead{\parbox[t]{0.58\textwidth}{\raggedright\small\sffamily\color{brandblue}\textbf{T11 --- Package, Moduli e Regole di Visibilità}}}
\rhead{\small\sffamily\color{gray}Programmazione II -- UniVR (2025/2026)}
\cfoot{\small\sffamily Pagina \thepage\ di \pageref{LastPage}}
\renewcommand{\headrulewidth}{0.5pt}
\renewcommand{\headrule}{\hbox to\headwidth{\color{brandblue!70!black}\leaders\hrule height \headrulewidth\hfill}}

\usepackage{titlesec}
\titleformat{\section}{\Large\bfseries\sffamily\color{brandblue}}{\thesection}{0.8em}{}[\titlerule]
\titleformat{\subsection}{\large\bfseries\sffamily\color{accentblue}}{\thesubsection}{0.8em}{}
\titleformat{\subsubsection}{\normalsize\bfseries\sffamily\color{brandblue!85!black}}{\thesubsubsection}{0.8em}{}

\usepackage{needspace}
\usepackage{fancyvrb}
\DefineVerbatimEnvironment{asciibox}{Verbatim}{
  fontsize=\fontsize{7.5pt}{9.2pt}\selectfont,
  frame=single,
  rulecolor=\color{codeframe},
  fillcolor=\color{codebg},
  framesep=2.5mm
}
\usepackage{listings}
\lstset{
  backgroundcolor=\color{codebg},
  frame=single,
  rulecolor=\color{codeframe},
  basicstyle=\ttfamily\footnotesize,
  keywordstyle=\color{javakeyword}\bfseries,
  stringstyle=\color{javastring},
  commentstyle=\color{javacomment}\itshape,
  numberstyle=\tiny\color{gray},
  numbers=left,
  numbersep=7pt,
  stepnumber=1,
  breaklines=true,
  breakatwhitespace=true,
  showstringspaces=false,
  tabsize=4,
  xleftmargin=12pt,
  framexleftmargin=12pt
}

\newcommand{\passthrough}[1]{#1}
\providecommand{\tightlist}{\setlength{\itemsep}{0pt}\setlength{\parskip}{0pt}}

\usepackage{tcolorbox}
\tcbuselibrary{skins,breakable}
\newtcolorbox{studynote}[1][Nota]{
  enhanced,
  breakable,
  colback=notebg,
  colframe=noteborder,
  fonttitle=\bfseries\sffamily\small,
  coltitle=white,
  title={#1},
  arc=2mm,
  boxrule=0.8pt,
  left=9pt, right=9pt, top=7pt, bottom=7pt
}

\newtcolorbox{studywarning}[1][Attenzione]{
  enhanced,
  breakable,
  colback=warnbg,
  colframe=warnborder,
  fonttitle=\bfseries\sffamily\small,
  coltitle=white,
  title={#1},
  arc=2mm,
  boxrule=0.8pt,
  left=9pt, right=9pt, top=7pt, bottom=7pt
}

\newtcolorbox{studytip}[1][Suggerimento]{
  enhanced,
  breakable,
  colback=tipbg,
  colframe=tipborder,
  fonttitle=\bfseries\sffamily\small,
  coltitle=white,
  title={#1},
  arc=2mm,
  boxrule=0.8pt,
  left=9pt, right=9pt, top=7pt, bottom=7pt
}

\usepackage{booktabs}
\usepackage{longtable}
\usepackage{array}
\usepackage{calc}

\usepackage{hyperref}
\hypersetup{
  colorlinks=true,
  linkcolor=accentblue,
  urlcolor=accentblue,
  citecolor=brandblue,
  pdfborder={0 0 0}
}

\title{\Huge\bfseries\sffamily\color{brandblue} T11 --- Package, Moduli e Regole di Visibilità \\[0.3em] \large\sffamily Struttura Cartelle, Modificatori public/protected/default/private e MavenDate v2 con equals()}
\author{\textbf{Docente:} Prof. Michele Pasqua \quad---\quad \textbf{Appunti Revisione:} Wally}
\date{A.A. 2025/2026 -- Universit\`a degli Studi di Verona}

\begin{document}
\maketitle
\thispagestyle{fancy}
\vspace{-1.2em}
\noindent\textcolor{brandblue}{\rule{\textwidth}{0.8pt}}
\vspace{1.2em}

In questo blocco analizziamo la modularità avanzata offerta dai package,
le regole di risoluzione dei nomi e i livelli di visibilità dei membri
(\href{file:///home/wally/Documenti/GitHub/Programmazione-II/25-26/Teoria-e-Esercitazioni/Slides-20260902/T11\%20-\%20Packages\%20and\%20Class\%20Visibility.pdf}{T11
- Packages and Class Visibility.pdf}). Sul piano pratico, in
\href{file:///home/wally/Documenti/GitHub/Programmazione-II/25-26/Teoria-e-Esercitazioni/Codice-20260902/Lezione11/MavenDate}{\passthrough{\lstinline!Lezione11/MavenDate!}}
assistiamo a una grande svolta architetturale: 1. Nascono le sottoclassi
specializzate
\textbf{\href{file:///home/wally/Documenti/GitHub/Programmazione-II/25-26/Teoria-e-Esercitazioni/Codice-20260902/Lezione11/MavenDate/src/main/java/it/oop/core/ItalianDate.java}{\passthrough{\lstinline!ItalianDate!}}}
e
\textbf{\href{file:///home/wally/Documenti/GitHub/Programmazione-II/25-26/Teoria-e-Esercitazioni/Codice-20260902/Lezione11/MavenDate/src/main/java/it/oop/core/AmericanDate.java}{\passthrough{\lstinline!AmericanDate!}}}
tramite estensione (\passthrough{\lstinline!extends Date!}). 2. Viene
introdotto il modificatore \textbf{\passthrough{\lstinline!protected!}}.
3. Viene implementato l'algoritmo canonico per l'overriding del metodo
\textbf{\passthrough{\lstinline!equals(Object)!}}. 4. Viene risolto il
bug storico dei getter \passthrough{\lstinline!getMonth()!} e
\passthrough{\lstinline!getYear()!}.

\begin{center}\rule{0.5\linewidth}{0.5pt}\end{center}

\subsubsection{1. Teoria: Concetti Teorici dalle Slide (Slide
T11)}\label{teoria-concetti-teorici-dalle-slide-slide-t11}

\paragraph{A. Modularità, Gerarchia e Convenzioni dei Package (Slide
1--11)}\label{a.-modularituxe0-gerarchia-e-convenzioni-dei-package-slide-111}

\begin{itemize}
\tightlist
\item
  \textbf{Definizione di Package}:

  \begin{itemize}
  \tightlist
  \item
    Un package è un raggruppamento logico di classi correlate che scala
    il concetto di incapsulamento a livello di libreria.
  \item
    Definisce un \textbf{namespace isolato}: classi con lo stesso nome
    possono coesistere senza conflitti se collocate in package
    differenti (es. \passthrough{\lstinline!it.oop.ui.Date!} vs
    \passthrough{\lstinline!it.oop.core.Date!}).
  \end{itemize}
\item
  \textbf{Convenzione di Naming Standard}:

  \begin{itemize}
  \tightlist
  \item
    Si usa il nome di dominio Internet invertito in caratteri minuscoli:
    \passthrough{\lstinline!it.univr.mypackage!},
    \passthrough{\lstinline!it.oop.core!}.
  \end{itemize}
\item
  \textbf{Dichiarazione e Corrispondenza col Filesystem}:

  \begin{itemize}
  \tightlist
  \item
    La direttiva \passthrough{\lstinline!package nomepkg;!} deve essere
    \textbf{la prima istruzione non di commento} all'inizio del file
    \passthrough{\lstinline!.java!}.
  \item
    Java impone una corrispondenza 1-a-1 tra la gerarchia dei package e
    la struttura delle cartelle su disco: la classe
    \passthrough{\lstinline!it.oop.core.Date!} deve trovarsi nel
    percorso \passthrough{\lstinline!it/oop/core/Date.java!}.
  \end{itemize}
\item
  \textbf{Sotto-Package e Wildcard}:

  \begin{itemize}
  \tightlist
  \item
    I package possono essere annidati (es.
    \passthrough{\lstinline!java.awt.event!} è un sottopackage di
    \passthrough{\lstinline!java.awt!}).
  \item
    Attenzione: \textbf{Regola d'esame}: la direttiva
    \passthrough{\lstinline!import pkg.*;!} importa solo le classi
    direttamente contenute in \passthrough{\lstinline!pkg!}. \textbf{Non
    importa le classi contenute nei suoi sotto-package} (per importare
    \passthrough{\lstinline!java.awt.event!} è obbligatorio scrivere
    \passthrough{\lstinline!import java.awt.event.*;!}).
  \end{itemize}
\item
  \textbf{Il Default Package (Unnamed Package)}:

  \begin{itemize}
  \tightlist
  \item
    Se un file \passthrough{\lstinline!.java!} omette la dichiarazione
    \passthrough{\lstinline!package!}, la classe appartiene al package
    predefinito senza nome.
  \item
    \textbf{Limitazione critica}: Le classi nel default package
    \textbf{non possono essere importate né utilizzate} da classi
    situate in package con nome. Nel software professionale e agli esami
    il suo uso è fortemente scoraggiato.
  \end{itemize}
\item
  \textbf{Compilazione e Packaging con Strumenti Standard e Maven}:

  \begin{itemize}
  \tightlist
  \item
    \passthrough{\lstinline!javac -d bin/ src/pkg/MyClass.java!}: il
    flag \passthrough{\lstinline!-d!} crea automaticamente l'albero
    delle sottocartelle nella directory di destinazione
    \passthrough{\lstinline!bin/!}.
  \item
    Con Maven: il layout \passthrough{\lstinline!src/main/java!} e
    \passthrough{\lstinline!src/test/java!} unito al
    \passthrough{\lstinline!pom.xml!} automatizza compilazione, test
    surefire e packaging in Fat JAR.
  \end{itemize}
\end{itemize}

\begin{center}\rule{0.5\linewidth}{0.5pt}\end{center}

\paragraph{B. Visibilità di Classe e Matrice di Accesso dei Membri
(Slide
12--19)}\label{b.-visibilituxe0-di-classe-e-matrice-di-accesso-dei-membri-slide-1219}

\begin{itemize}
\tightlist
\item
  \textbf{Visibilità a Livello di Classe (Top-Level Classes)}:

  \begin{itemize}
  \tightlist
  \item
    \passthrough{\lstinline!public class Clazz!}: accessibile e
    istanziabile da qualsiasi package del Classpath.
  \item
    \passthrough{\lstinline!class Clazz!} (\emph{Package-Private} /
    default): visibile e utilizzabile \textbf{esclusivamente all'interno
    dello stesso package}. Non può essere usata all'esterno, anche se i
    suoi metodi interni sono dichiarati
    \passthrough{\lstinline!public!}.
  \item
    \emph{Nota}: le classi di primo livello non possono essere
    dichiarate \passthrough{\lstinline!private!} né
    \passthrough{\lstinline!protected!} (questi modificatori sono
    ammessi solo per classi annidate/interne).
  \end{itemize}
\item
  \textbf{La Matrice dei Modificatori di Accesso (Fields e Methods)}:
\end{itemize}

\begin{longtable}[]{@{}
  >{\raggedright\arraybackslash}p{(\linewidth - 8\tabcolsep) * \real{0.1667}}
  >{\centering\arraybackslash}p{(\linewidth - 8\tabcolsep) * \real{0.2083}}
  >{\centering\arraybackslash}p{(\linewidth - 8\tabcolsep) * \real{0.2083}}
  >{\centering\arraybackslash}p{(\linewidth - 8\tabcolsep) * \real{0.2083}}
  >{\centering\arraybackslash}p{(\linewidth - 8\tabcolsep) * \real{0.2083}}@{}}
\toprule\noalign{}
\begin{minipage}[b]{\linewidth}\raggedright
Modificatore
\end{minipage} & \begin{minipage}[b]{\linewidth}\centering
Stessa Classe
\end{minipage} & \begin{minipage}[b]{\linewidth}\centering
Stesso Package
\end{minipage} & \begin{minipage}[b]{\linewidth}\centering
Sottoclasse (Altro Package)
\end{minipage} & \begin{minipage}[b]{\linewidth}\centering
Mondo Esterno
\end{minipage} \\
\midrule\noalign{}
\endhead
\bottomrule\noalign{}
\endlastfoot
\textbf{\passthrough{\lstinline!private!}} & {[}OK{]} & {[}NO{]} &
{[}NO{]} & {[}NO{]} \\
\textbf{Default} \emph{(package-private)} & {[}OK{]} & {[}OK{]} &
{[}NO{]} & {[}NO{]} \\
\textbf{\passthrough{\lstinline!protected!}} & {[}OK{]} & {[}OK{]} &
{[}OK{]} & {[}NO{]} \\
\textbf{\passthrough{\lstinline!public!}} & {[}OK{]} & {[}OK{]} &
{[}OK{]} & {[}OK{]} \\
\end{longtable}

\begin{itemize}
\tightlist
\item
  \textbf{Risoluzione dei Nomi}:

  \begin{itemize}
  \tightlist
  \item
    Se una classe importa due tipi omonimi o usa una classe locale
    avente lo stesso nome di una classe importata, per disambiguare è
    obbligatorio usare il \textbf{Fully Qualified Name (FQN)}:
    \passthrough{\lstinline!it.oop.ui.Date d = new it.oop.ui.Date(12345);!}
  \end{itemize}
\end{itemize}

\begin{center}\rule{0.5\linewidth}{0.5pt}\end{center}

\subsubsection{\texorpdfstring{2. Codice: Evoluzione del Codice:
\href{file:///home/wally/Documenti/GitHub/Programmazione-II/25-26/Teoria-e-Esercitazioni/Codice-20260902/Lezione11/MavenDate}{\texttt{Lezione11/MavenDate}}}{2. Codice: Evoluzione del Codice: Lezione11/MavenDate}}\label{codice-evoluzione-del-codice-lezione11mavendate}

In \passthrough{\lstinline!Lezione11!}, il codice di
\passthrough{\lstinline!MavenDate!} compie un balzo in avanti con
l'introduzione di gerarchie di classi e l'overriding:

\paragraph{\texorpdfstring{1. Generalizzazione in
\href{file:///home/wally/Documenti/GitHub/Programmazione-II/25-26/Teoria-e-Esercitazioni/Codice-20260902/Lezione11/MavenDate/src/main/java/it/oop/core/Date.java}{\texttt{Date.java}}}{1. Generalizzazione in Date.java}}\label{generalizzazione-in-date.java}

\begin{lstlisting}[language=Java]
package it.oop.core;

public class Date {
    // 1. Campo 'protected': accessibile direttamente dalle sottoclassi
    protected final int day;
    // 2. Campi 'private': incapsulati e accessibili solo tramite getter
    private final int month;
    private final int year;

    public Date(int day, int month, int year) {
        this.day = day;
        this.month = month;
        this.year = year;
        verify();
    }

    // 3. I GETTER SONO STATI CORRETTI!
    public int getDay() { return day; }
    public int getMonth() { return month; } // Non restituisce più day!
    public int getYear() { return year; }   // Non restituisce più day!

    // 4. Implementazione canonica di equals(Object)
    @Override
    public boolean equals(Object other) {
        if (other == null) return false;              // Controllo null-safety
        if (this == other) return true;               // Controllo identità (stesso indirizzo)
        if (!(other instanceof Date)) return false;   // Controllo di tipo/ereditarietà
        Date otherAsDate = (Date) other;              // Downcast sicuro
        return day == otherAsDate.getDay()            // Confronto dei campi
                && month == otherAsDate.getMonth()
                && year == otherAsDate.getYear();
    }
}
\end{lstlisting}

\paragraph{\texorpdfstring{Analisi: I 5 Passi Fondamentali
dell'Algoritmo
\texttt{equals(Object)}:}{Analisi: I 5 Passi Fondamentali dell'Algoritmo equals(Object):}}\label{analisi-i-5-passi-fondamentali-dellalgoritmo-equalsobject}

\begin{enumerate}
\def\labelenumi{\arabic{enumi}.}
\tightlist
\item
  \passthrough{\lstinline!if (other == null) return false;!}: garantisce
  che \passthrough{\lstinline!d.equals(null)!} ritorni
  \passthrough{\lstinline!false!} senza lanciare
  \passthrough{\lstinline!NullPointerException!}.
\item
  \passthrough{\lstinline!if (this == other) return true;!}:
  ottimizzazione immediata (riflessività). Se sono lo stesso oggetto
  nello Heap, sono identici.
\item
  \passthrough{\lstinline"if (!(other instanceof Date)) return false;"}:
  verifica se l'oggetto passato è compatibile con la gerarchia
  \passthrough{\lstinline!Date!}.
\item
  \passthrough{\lstinline!Date otherAsDate = (Date) other;!}: downcast
  necessario perché la firma di \passthrough{\lstinline!equals!} accetta
  un generico \passthrough{\lstinline!Object!}.
\item
  Confronto logico dei campi che definiscono lo stato:
  \passthrough{\lstinline!day!}, \passthrough{\lstinline!month!},
  \passthrough{\lstinline!year!}.
\end{enumerate}

\begin{center}\rule{0.5\linewidth}{0.5pt}\end{center}

\paragraph{\texorpdfstring{2. Le Nuove Sottoclassi:
\href{file:///home/wally/Documenti/GitHub/Programmazione-II/25-26/Teoria-e-Esercitazioni/Codice-20260902/Lezione11/MavenDate/src/main/java/it/oop/core/ItalianDate.java}{\texttt{ItalianDate.java}}
e
\href{file:///home/wally/Documenti/GitHub/Programmazione-II/25-26/Teoria-e-Esercitazioni/Codice-20260902/Lezione11/MavenDate/src/main/java/it/oop/core/AmericanDate.java}{\texttt{AmericanDate.java}}}{2. Le Nuove Sottoclassi: ItalianDate.java e AmericanDate.java}}\label{le-nuove-sottoclassi-italiandate.java-e-americandate.java}

\begin{lstlisting}[language=Java]
package it.oop.core;

public class ItalianDate extends Date {
    private static final String FORMAT = "dd/mm/yyyy";
    private static final String[] MONTHS = { "gennaio", "febbraio", ... };

    // Invocazione del costruttore genitore con 'super'
    public ItalianDate(int day, int month, int year) {
        super(day, month, year);
    }

    public String printFormat() { return FORMAT; }

    public String prettyPrint() {
        // 'day' è accessibile direttamente perché 'protected'!
        // 'getYear()' è invocato tramite metodo perché 'year' è 'private' in Date!
        return day + " " + getMonthAsString() + " " + getYear();
    }

    @Override
    public String toString() {
        return day + "/" + getMonth() + "/" + getYear();
    }
}
\end{lstlisting}

E simmetricamente per \passthrough{\lstinline!AmericanDate!}:

\begin{lstlisting}[language=Java]
public class AmericanDate extends Date {
    private static final String FORMAT = "mm/dd/yyyy";

    public AmericanDate(int day, int month, int year) {
        super(day, month, year);
    }

    @Override
    public String toString() {
        // Mese prima del giorno!
        return getMonth() + "/" + getDay() + "/" + getYear();
    }
}
\end{lstlisting}

\begin{center}\rule{0.5\linewidth}{0.5pt}\end{center}

\paragraph{\texorpdfstring{3. Polimorfismo, Upcasting e Downcasting in
\href{file:///home/wally/Documenti/GitHub/Programmazione-II/25-26/Teoria-e-Esercitazioni/Codice-20260902/Lezione11/MavenDate/src/main/java/it/oop/ui/MainDate.java}{\texttt{MainDate.java}}}{3. Polimorfismo, Upcasting e Downcasting in MainDate.java}}\label{polimorfismo-upcasting-e-downcasting-in-maindate.java}

Analizziamo i comportamenti a runtime testati dal docente:

\begin{lstlisting}[language=Java]
Date date = new Date(3, 11, 2025);
ItalianDate itDate = new ItalianDate(3, 11, 2025);
AmericanDate usDate = new AmericanDate(3, 11, 2025);

// 1. Confronti con equals():
System.out.println("date ?= itDate: " + date.equals(itDate));     // Stampa TRUE!
System.out.println("itDate ?= usDate: " + itDate.equals(usDate)); // Stampa TRUE!
\end{lstlisting}

\emph{Perché stampano \passthrough{\lstinline!true!}?} Perché sia
\passthrough{\lstinline!itDate!} che \passthrough{\lstinline!usDate!}
estendono \passthrough{\lstinline!Date!}: l'operatore
\passthrough{\lstinline!instanceof Date!} è soddisfatto, e i campi
(giorno 3, mese 11, anno 2025) coincidono perfettamente!

\begin{lstlisting}[language=Java]
// 2. Upcasting (Polimorfismo di Inclusione):
Date d = new ItalianDate(1, 1, 1970); // LECITO: ItalianDate È UN Date!

// 3. Tipi Incompatibili tra rami fratelli:
// AmericanDate ad = new ItalianDate(1, 1, 1970); // COMPILE-TIME ERROR!

// 4. Downcasting esplicito lecito:
ItalianDate id = (ItalianDate) d; // LECITO a runtime: l'oggetto nello Heap è davvero ItalianDate!
System.out.println(id.printFormat()); // Stampa dd/mm/yyyy

// 5. Downcasting illegale a runtime:
// ItalianDate idFalso = (ItalianDate) date; // CRASH RUNTIME: ClassCastException!
// (perché 'date' è un Date generico, non contiene la specializzazione ItalianDate)
\end{lstlisting}

\begin{center}\rule{0.5\linewidth}{0.5pt}\end{center}

\subsubsection{3. Ciclo: Corrispondenza con i Tuoi Moduli
Workspace}\label{ciclo-corrispondenza-con-i-tuoi-moduli-workspace}

Nel tuo repository
\href{file:///home/wally/Documenti/GitHub/Programmazione-II/25-26/esercizi-wally-25-26}{\passthrough{\lstinline!esercizi-wally-25-26!}}:
- \textbf{\passthrough{\lstinline!es13!}}:
\href{file:///home/wally/Documenti/GitHub/Programmazione-II/25-26/esercizi-wally-25-26/es13/src/main/java/es13/model/Calculator.java}{\passthrough{\lstinline!Calculator.java!}}
e
\href{file:///home/wally/Documenti/GitHub/Programmazione-II/25-26/esercizi-wally-25-26/es13/src/test/java/es13/TestCalculator.java}{\passthrough{\lstinline!TestCalculator.java!}}
applicano la separazione \passthrough{\lstinline!es13.model!} e
\passthrough{\lstinline!es13!} illustrata a lezione (Slide T11,
pp.~3--14). - \textbf{\passthrough{\lstinline!es14!}}: implementa
esattamente questa architettura a oggetti multi-package
(\passthrough{\lstinline!it.oop.core!} e
\passthrough{\lstinline!it.oop.ui!}) con
\href{file:///home/wally/Documenti/GitHub/Programmazione-II/25-26/esercizi-wally-25-26/es14/src/main/java/it/oop/core/ItalianDate.java}{\passthrough{\lstinline!ItalianDate!}},
\href{file:///home/wally/Documenti/GitHub/Programmazione-II/25-26/esercizi-wally-25-26/es14/src/main/java/it/oop/core/AmericanDate.java}{\passthrough{\lstinline!AmericanDate!}},
\href{file:///home/wally/Documenti/GitHub/Programmazione-II/25-26/esercizi-wally-25-26/es14/src/main/java/it/oop/core/Date.java}{\passthrough{\lstinline!Date!}}
e i relativi test in
\href{file:///home/wally/Documenti/GitHub/Programmazione-II/25-26/esercizi-wally-25-26/es14/src/test/java/it/oop/core/TestItalianDate.java}{\passthrough{\lstinline!src/test/java/it/oop/core/TestItalianDate.java!}}.

\begin{center}\rule{0.5\linewidth}{0.5pt}\end{center}

\begin{studynote}[Nota di Studio] Con la \textbf{Lezione 11} abbiamo visto l'introduzione
dell'ereditarietà (\passthrough{\lstinline!extends!}), la visibilità
\passthrough{\lstinline!protected!}, l'overriding di
\passthrough{\lstinline!equals()!} e le regole di cast tra classi.

Il prossimo blocco è la \textbf{Lezione 12 / Slide T12:
\emph{Inheritance and Polymorphism}}: - Teoria approfondita
sull'ereditarietà: riuso, estensione e specializzazione. - La classe
radice universale \textbf{\passthrough{\lstinline!java.lang.Object!}} e
i suoi metodi (\passthrough{\lstinline!toString!},
\passthrough{\lstinline!equals!}, \passthrough{\lstinline!hashCode!},
\passthrough{\lstinline!getClass!}, \passthrough{\lstinline!clone!}). -
Il costrutto \passthrough{\lstinline!super!}: invocazione di costruttori
(\passthrough{\lstinline!super(...)!}) e invocazione di metodi della
superclasse (\passthrough{\lstinline!super.metodo()!}). -
\textbf{Polimorfismo e Dynamic Method Dispatch (Late Binding)}: come la
JVM decide a runtime quale metodo eseguire tramite la \emph{vtable}. -
Regole di overriding vs overloading e l'annotazione
\passthrough{\lstinline!@Override!}. - Evoluzione del codice in
\passthrough{\lstinline!Lezione12!}: - La classe
\passthrough{\lstinline!Date!} diventa la superclasse polimorfica di
riferimento. - Aggiunta del metodo
\passthrough{\lstinline!isLeapYear(int year)!} per la validazione
completa dei bisestili gregoriani!
\end{studynote}

\textbf{Dammi conferma per aprire la Lezione 12 / T12 e proseguire!}


\end{document}
